\phantomsection\addcontentsline{toc}{section}{Marine Renewable Energy}
\ECchapter{25}{Marine Renewable Energy}{How can oceans supply dependable renewable power?}{ECBlue}

\ECObjectives{
\begin{itemize}
\item Compare tidal, wave and offshore energy technologies.
\item Explain the mechanical, environmental and maintenance constraints.
\item Estimate power, annual energy yield and capacity factor.
\end{itemize}}

\begin{multicols}{2}
\begin{engineeringnews}
Oceans contain major renewable resources. Offshore wind is already deployed at commercial scale,
whereas tidal-stream and wave-energy devices are still developing. Marine systems benefit from strong,
partly predictable resources but face demanding installation, grid-connection and maintenance conditions.
\end{engineeringnews}

\begin{ecarticle}
Marine renewable energy includes offshore wind, tidal barrages, tidal-stream turbines and wave-energy
converters. Each technology extracts energy from a different physical phenomenon and therefore requires
a different structure, control strategy and maintenance plan.

Tidal-stream turbines operate like underwater wind turbines. Because water is much denser than air,
even a moderate current can exert large forces. The available kinetic power is
\[
P_{\mathrm{avail}}=\frac{1}{2}\rho A v^3,
\]
where $\rho$ is water density, $A$ the swept area and $v$ the current speed. The electrical output is
lower because rotor, generator and converter efficiencies are limited. Tides are predictable, but
suitable sites are limited and the equipment must resist strong cyclic loads.

Wave-energy converters use relative motion between waves and a floating or fixed structure. Some use
oscillating water columns; others use hinged bodies or hydraulic power take-off systems. Wave direction,
period and height vary continuously, so the controller must adapt while limiting excessive motion.
Survival during an extreme storm may be more demanding than normal energy production.

Salt water accelerates corrosion and biological growth. Barnacles and algae modify surface properties
and increase drag. Seals, cables and bearings must operate for long periods with limited access.
Retrieving a device may require a specialised vessel and a calm weather window, so reliability and
maintainability strongly influence the levelised cost of energy.

Engineers must also consider navigation, fishing, underwater noise and marine ecosystems. The best
design is therefore not simply the device with the highest peak power, but the system that delivers
dependable energy over its complete life cycle.
\end{ecarticle}

\begin{tcolorbox}[title=\sffamily\bfseries Engineering Insight,colback=yellow!10,colframe=ECOrange]
Marine devices are often designed by two operating cases: efficient energy capture in normal conditions
and structural survival during rare extreme events.
\end{tcolorbox}

\begin{engineeringfocus}
\textbf{Tidal-stream conversion chain}
\begin{center}
\resizebox{.96\linewidth}{!}{%
\begin{tikzpicture}[font=\scriptsize,>=Latex,node distance=3mm]
\node[draw=ECBlue,fill=ECLightBlue,rounded corners] (flow) {tidal current};
\node[draw=ECGreen,fill=ECGreen!10,rounded corners,right=of flow] (rotor) {rotor};
\node[draw=ECPurple,fill=ECPurple!10,rounded corners,right=of rotor] (gen) {generator};
\node[draw=ECOrange,fill=ECOrange!10,rounded corners,right=of gen] (conv) {converter};
\node[draw=ECRed,fill=ECRed!8,rounded corners,below=of gen] (grid) {subsea cable and grid};
\draw[->,thick] (flow)--(rotor); \draw[->,thick] (rotor)--(gen); \draw[->,thick] (gen)--(conv); \draw[->,thick] (conv)--(grid);
\end{tikzpicture}%
}
\end{center}
Each conversion stage introduces losses and a possible failure mode that must be inspected or monitored.
\end{engineeringfocus}

\begin{keyfigures}
\begin{tabularx}{\linewidth}{@{}Xr@{}}
Predictable resource & tides\\
Main material threat & corrosion\\
Common surface issue & biofouling\\
Major cost driver & offshore maintenance\\
Grid connection & subsea cable
\end{tabularx}
\end{keyfigures}

\begin{tcolorbox}[title=\sffamily\bfseries Illustrative capacity factors,colback=white,colframe=ECBlue]
\begin{center}
\begin{tikzpicture}
\begin{axis}[xbar,width=.96\linewidth,height=48mm,xlabel={Capacity factor (\%)},xmin=0,xmax=55,symbolic y coords={Tidal barrage,Offshore wind,Wave energy,Tidal stream},ytick=data,grid=major,ticklabel style={font=\scriptsize},nodes near coords]
\addplot table[x=factor,y=technology,col sep=comma]{data/EC25/capacity.csv};
\end{axis}
\end{tikzpicture}
\end{center}
\footnotesize Illustrative values; real performance depends strongly on the site, availability and technology maturity.
\end{tcolorbox}

\begin{vocabulary}
\begin{tabularx}{\linewidth}{@{}lX@{}}
tidal stream & courant de marée\\
wave period & période de vague\\
power take-off & système de conversion d'énergie\\
cyclic load & charge cyclique\\
corrosion & corrosion\\
biofouling & encrassement biologique\\
subsea cable & câble sous-marin\\
capacity factor & facteur de charge\\
mooring & amarrage\\
storm survival & tenue en tempête
\end{tabularx}
\end{vocabulary}

\begin{questions}
\item Why can relatively slow seawater flows produce large forces?
\item What makes wave-energy control difficult?
\item Why does maintenance strongly influence marine-energy costs?
\item Give two environmental impacts that must be monitored.
\item Why is peak power insufficient for comparing devices?
\end{questions}

\begin{applicationexercise}
A tidal turbine has a rotor diameter of $18\,\mathrm{m}$ in a current of $2.5\,\mathrm{m\,s^{-1}}$.
Take $\rho=1025\,\mathrm{kg\,m^{-3}}$, $C_p=0.40$ and an electrical efficiency of $90\%$.
\begin{enumerate}
\item Calculate the swept area.
\item Estimate the electrical power produced at this current speed.
\item With a capacity factor of $35\%$, estimate the annual energy output.
\item The device is unavailable for maintenance for 20 days per year. Recalculate the annual output.
\item Explain why the $v^3$ relationship makes site measurements critical.
\end{enumerate}
\end{applicationexercise}

\begin{engineeringchallenge}
Propose a nominal 1 MW tidal-stream installation. Estimate the rotor area for an assumed current speed,
select corrosion and biofouling protection, define condition-monitoring sensors and prepare a maintenance
strategy. Add two ecological indicators and one storm-survival verification test.
\end{engineeringchallenge}

\begin{speaking}
Should marine renewable projects be prioritised in protected coastal areas when they reduce carbon
emissions? Discuss biodiversity, energy security, local jobs and precaution.
\end{speaking}
\end{multicols}

\subsection*{Sources}
\ECsource{International Energy Agency -- Ocean Power}{https://www.iea.org/energy-system/renewables/ocean-power}
\ECsource{European Marine Energy Centre}{https://www.emec.org.uk/}